\documentclass[12pt, a4paper]{book}

% ─── Core packages ────────────────────────────────────────────────────────────
\usepackage[T1]{fontenc}
\usepackage[utf8]{inputenc}
\usepackage[charter, cal=cmcal]{mathdesign}
\usepackage[scaled=0.92]{helvet}
\usepackage{geometry}
\geometry{
  a4paper,
  top=30mm, bottom=30mm,
  inner=32mm, outer=25mm,
  marginparwidth=18mm
}

% ─── Mathematics ──────────────────────────────────────────────────────────────
\usepackage{amsmath, amsthm}

% ─── Chemistry ────────────────────────────────────────────────────────────────
\usepackage[version=4]{mhchem}

% ─── Graphics & colour ────────────────────────────────────────────────────────
\usepackage{graphicx}
\usepackage[dvipsnames,svgnames,table]{xcolor}
\usepackage{tikz}
\usetikzlibrary{arrows.meta, shapes.geometric, positioning,
                calc, backgrounds, patterns, decorations.pathmorphing}
\usepackage{pgfplots}
\pgfplotsset{compat=1.18}
\usepackage{tcolorbox}
\tcbuselibrary{skins, breakable, theorems}

% ─── Layout ───────────────────────────────────────────────────────────────────
\usepackage{titlesec}
\usepackage{fancyhdr}
\usepackage{enumitem}
\usepackage{booktabs}
\usepackage{array}
\usepackage{multicol}
\usepackage{marginnote}
\usepackage{mdframed}
\usepackage{wrapfig}
\usepackage{caption}
\usepackage{subcaption}
\usepackage{parskip}
\usepackage{siunitx}
\usepackage[colorlinks=true, linkcolor=BrickRed,
            urlcolor=NavyBlue, citecolor=ForestGreen]{hyperref}

% ══════════════════════════════════════════════════════════════════════════════
% COLOUR PALETTE
% ══════════════════════════════════════════════════════════════════════════════
\definecolor{JourneyColor}{HTML}{065F46}   % forest green — journey headings
\definecolor{ModuleColor}{HTML}{0F766E}    % deep teal — module headings
\definecolor{PuzzleColor}{HTML}{B45309}    % deep amber — puzzle boxes
\definecolor{DiscoveryColor}{HTML}{D97706} % golden amber — discovery boxes
\definecolor{WarningColor}{HTML}{E11D48}   % coral red — "model breaks" boxes
\definecolor{ProblemColor}{HTML}{0F766E}   % dark teal — problem sets
\definecolor{ReflectColor}{HTML}{047857}   % green — reflection boxes
\definecolor{HistoryColor}{HTML}{9A3412}   % brown — history notes
\definecolor{LightBg}{HTML}{F8FAFC}        % near-white background

% ══════════════════════════════════════════════════════════════════════════════
% TCOLORBOX ENVIRONMENTS
% ══════════════════════════════════════════════════════════════════════════════
\newtcolorbox{puzzlebox}[1][]{
  enhanced, breakable,
  colback=PuzzleColor!6, colframe=PuzzleColor,
  fonttitle=\bfseries\sffamily, title={The Puzzle},
  attach boxed title to top left={yshift=-2mm, xshift=6mm},
  boxed title style={colback=PuzzleColor, colframe=PuzzleColor},
  borderline west={3pt}{0pt}{PuzzleColor}, #1}

\newtcolorbox{discoverybox}[1][]{
  enhanced, breakable,
  colback=DiscoveryColor!8, colframe=DiscoveryColor,
  fonttitle=\bfseries\sffamily, title={Discovery},
  attach boxed title to top left={yshift=-2mm, xshift=6mm},
  boxed title style={colback=DiscoveryColor, colframe=DiscoveryColor},
  borderline west={3pt}{0pt}{DiscoveryColor}, #1}

\newtcolorbox{problembox}[1][]{
  enhanced, breakable,
  colback=ProblemColor!5, colframe=ProblemColor,
  fonttitle=\bfseries\sffamily, title={Problems},
  attach boxed title to top left={yshift=-2mm, xshift=6mm},
  boxed title style={colback=ProblemColor, colframe=ProblemColor},
  borderline west={3pt}{0pt}{ProblemColor}, #1}

\newtcolorbox{deepbox}[1][]{
  enhanced, breakable,
  colback=JourneyColor!5, colframe=JourneyColor,
  fonttitle=\bfseries\sffamily, title={Deep Problems},
  attach boxed title to top left={yshift=-2mm, xshift=6mm},
  boxed title style={colback=JourneyColor, colframe=JourneyColor},
  borderline west={4pt}{0pt}{JourneyColor}, #1}

\newtcolorbox{historybox}[1][]{
  enhanced, breakable,
  colback=HistoryColor!7, colframe=HistoryColor!60,
  fonttitle=\bfseries\sffamily\small, title={Historical Note},
  attach boxed title to top left={yshift=-2mm, xshift=6mm},
  boxed title style={colback=HistoryColor!80, colframe=HistoryColor}, #1}

\newtcolorbox{reflectionbox}[1][]{
  enhanced, breakable,
  colback=ReflectColor!6, colframe=ReflectColor,
  fonttitle=\bfseries\sffamily, title={Reflection},
  attach boxed title to top left={yshift=-2mm, xshift=6mm},
  boxed title style={colback=ReflectColor, colframe=ReflectColor},
  borderline west={3pt}{0pt}{ReflectColor}, #1}

\newtcolorbox{modelbreaks}[1][]{
  enhanced, breakable,
  colback=WarningColor!6, colframe=WarningColor,
  fonttitle=\bfseries\sffamily, title={Where This Model Breaks},
  attach boxed title to top left={yshift=-2mm, xshift=6mm},
  boxed title style={colback=WarningColor, colframe=WarningColor},
  borderline west={4pt}{0pt}{WarningColor}, #1}

\newtcolorbox{namedconcept}[2][]{
  enhanced,
  colback=ModuleColor!6, colframe=ModuleColor,
  fonttitle=\bfseries\sffamily\large, title={#2},
  attach boxed title to top left={yshift=-2mm, xshift=6mm},
  boxed title style={colback=ModuleColor, colframe=ModuleColor}, #1}

% ══════════════════════════════════════════════════════════════════════════════
% THEOREM ENVIRONMENTS
% ══════════════════════════════════════════════════════════════════════════════
\theoremstyle{definition}
\newtheorem{observation}{Observation}[chapter]
\newtheorem{question}{Question}[chapter]
\newtheorem{experiment}{Experiment}[chapter]

% ══════════════════════════════════════════════════════════════════════════════
% SECTION FORMATTING
% ══════════════════════════════════════════════════════════════════════════════
\titleformat{\chapter}[display]
  {\normalfont\sffamily\huge\bfseries\color{JourneyColor}}
  {\textcolor{JourneyColor!60}{\chaptertitlename\ \thechapter}}
  {10pt}{\Huge}[\vspace{2pt}\textcolor{JourneyColor}{\titlerule[1.5pt]}]

\titleformat{\section}
  {\normalfont\sffamily\Large\bfseries\color{ModuleColor}}
  {\thesection}{1em}{}[\textcolor{ModuleColor!40}{\titlerule[0.6pt]}]

\titleformat{\subsection}
  {\normalfont\sffamily\large\bfseries\color{JourneyColor!80}}
  {\thesubsection}{1em}{}

\titleformat{\subsubsection}
  {\normalfont\sffamily\normalsize\bfseries}
  {\thesubsubsection}{1em}{}

% ══════════════════════════════════════════════════════════════════════════════
% HEADER / FOOTER
% ══════════════════════════════════════════════════════════════════════════════
\pagestyle{fancy}
\fancyhf{}
\fancyhead[LE]{\small\sffamily\textcolor{JourneyColor}{\leftmark}}
\fancyhead[RO]{\small\sffamily\textcolor{ModuleColor}{\rightmark}}
\fancyfoot[C]{\small\thepage}
\renewcommand{\headrulewidth}{0.4pt}

% ══════════════════════════════════════════════════════════════════════════════
% CUSTOM COMMANDS
% ══════════════════════════════════════════════════════════════════════════════
\newcommand{\chem}[1]{\ce{#1}}
\newcommand{\important}[1]{\textcolor{ModuleColor}{\textbf{#1}}}
\newcommand{\newterm}[1]{\textit{\textcolor{DiscoveryColor}{#1}}}
\newcommand{\sidenote}[1]{%
  \marginnote{\footnotesize\sffamily\textcolor{JourneyColor!70}{#1}}}

% ══════════════════════════════════════════════════════════════════════════════
% DOCUMENT
% ══════════════════════════════════════════════════════════════════════════════
\begin{document}
\setcounter{chapter}{0}
\chapter[Journey 1 --- What Is Stuff?]
        {Journey One\\[8pt]
         \textcolor{ModuleColor}{\large What Is Stuff?\ \ Matter and Its Patterns}}
\setcounter{section}{2}

% ══════════════════════════════════════════════════════════════════════════════
%  MODULE 1.3 — ELEMENTS, LAVOISIER'S TABLE, AND THE FIXED-RATIO MYSTERY
% ══════════════════════════════════════════════════════════════════════════════
\section{Module 1.3 --- Elements, Fixed Ratios, and the Demand for Atoms}
\label{sec:mod1.3}

\subsection*{Motivation}

By the end of Module~1.2 we had two powerful tools:
\begin{itemize}[leftmargin=2em, nosep]
  \item A precise definition of an \emph{element}
        (a substance that resists all chemical decomposition), and
  \item The law of conservation of mass
        (total mass is unchanged in any reaction).
\end{itemize}

\vspace{0.4em}

But something deeper and stranger was hiding in Lavoisier's data,
waiting to be noticed.

When hydrogen and oxygen combine to form water,
the ratio of their masses is always
\SI{1}{\gram} of hydrogen to \SI{8}{\gram} of oxygen.
Not sometimes. Not approximately. \emph{Always}, regardless
of source, temperature, or apparatus.

When carbon burns to form carbon dioxide,
the ratio is always
\SI{3}{\gram} of carbon to \SI{8}{\gram} of oxygen.

\begin{center}
\begin{tcolorbox}[
  enhanced, colback=LightBg, colframe=JourneyColor!40,
  width=0.86\textwidth, arc=3pt, fontupper=\sffamily
]
\centering
Why should elements always combine in \emph{fixed mass ratios}?\\[6pt]
\small
If matter were continuous --- like water poured into water ---\\
you could combine any amounts in any ratio.\\[4pt]
\normalsize
The \emph{fixed ratios} are a clue.\\
They are whispering something about the structure of matter.
\end{tcolorbox}
\end{center}

\vspace{0.5em}

In this module we will follow the data until it forces upon us
a single unavoidable conclusion:
matter must be made of \emph{discrete, indivisible particles}
--- atoms --- and each element has its own characteristic particle.

% ─── Section: Lavoisier's Table of Elements ───────────────────────────────────
\subsection{Lavoisier's Table of Elements: The First Inventory}

Before we reach fixed ratios, we need to know how many
elements there were to work with.

Lavoisier compiled the first modern list of elements in 1789,
in his textbook \textit{Traité Élémentaire de Chimie}.
His criterion: an element is any substance that resists
decomposition by any known chemical method.

He listed 33 substances as elements.
Most were correct.
Some were wrong in illuminating ways.

\begin{center}
\begin{tikzpicture}[
  every node/.style={font=\sffamily\footnotesize},
  box/.style={rounded corners=3pt, minimum height=0.7cm,
              minimum width=3.2cm, align=center}
]
  % Title
  \node[font=\sffamily\small\bfseries, color=JourneyColor]
        at (0, 2.0) {Lavoisier's 1789 Element List (selected)};

  % Three columns
  % Column 1: Correct elements
  \node[box, draw=ModuleColor!70, fill=ModuleColor!8,
        minimum width=4cm]
        at (-5.0, 0.8) {\textbf{Correct elements}};
  \foreach \elem/\y in {
    {Oxygen (\ce{O})}/0.0,
    {Hydrogen (\ce{H})}/-0.7,
    {Nitrogen (\ce{N})}/-1.4,
    {Sulfur (\ce{S})}/-2.1,
    {Carbon (\ce{C})}/-2.8,
    {Iron (\ce{Fe})}/-3.5,
    {Copper (\ce{Cu})}/-4.2,
    {Gold (\ce{Au})}/-4.9}
  {
    \node[box, draw=ModuleColor!40, fill=ModuleColor!5,
          minimum width=4cm] at (-5.0, \y) {\elem};
  }

  % Column 2: Wrong entries
  \node[box, draw=WarningColor!70, fill=WarningColor!8,
        minimum width=4cm]
        at (0.2, 0.8) {\textbf{Wrongly listed as elements}};
  \foreach \elem/\y in {
    {Light}/0.0,
    {Caloric (heat)}/-0.7,
    {Lime (\ce{CaO})}/-1.4,
    {Magnesia (\ce{MgO})}/-2.1,
    {Barytes (\ce{BaSO4})}/-2.8}
  {
    \node[box, draw=WarningColor!40, fill=WarningColor!5,
          minimum width=4cm] at (0.2, \y) {\elem};
  }

  % Column 3: Unknown at the time
  \node[box, draw=DiscoveryColor!70, fill=DiscoveryColor!8,
        minimum width=4cm]
        at (5.4, 0.8) {\textbf{Unknown in 1789}};
  \foreach \elem/\y in {
    {Sodium (\ce{Na})}/ 0.0,
    {Potassium (\ce{K})}/-0.7,
    {Calcium (\ce{Ca})}/-1.4,
    {Chlorine (\ce{Cl})}/-2.1,
    {Iodine (\ce{I})}/-2.8,
    {All noble gases}/-3.5}
  {
    \node[box, draw=DiscoveryColor!40, fill=DiscoveryColor!5,
          minimum width=4cm] at (5.4, \y) {\elem};
  }
\end{tikzpicture}
\end{center}

\vspace{0.5em}

The errors are instructive.
Lavoisier listed \emph{light} and \emph{caloric} (heat)
as elements because neither could be decomposed.
This shows the danger of a purely operational definition:
if your definition of ``element'' is
``what cannot be decomposed,''
you may inadvertently include things that are not matter at all.

The wrong entries were corrected over the following decades
as better instruments were invented.
Lime (\ce{CaO}) was decomposed into calcium and oxygen in 1808
by Humphry Davy using electrolysis.

\begin{historybox}
\textbf{Humphry Davy and the Electrolysis Revolution (1807--1808)}

Humphry Davy was a showman and a brilliant experimentalist.
In 1807 he passed an electric current through molten potash
(\ce{K2O}) and obtained tiny globules of a new metal
that burst into violet flame on contact with water.
He named it \emph{potassium}.

Days later, he repeated the process with molten soda (\ce{Na2O})
and obtained \emph{sodium} --- another metal that reacts violently
with water.

In 1808 he decomposed lime (\ce{CaO}), magnesia (\ce{MgO}),
and baryta (\ce{BaO}), revealing calcium, magnesium, and barium.

In a few months, Davy doubled the number of known elements.
His tool was electricity --- a force powerful enough to break
bonds that no chemical reagent could touch.

This episode reveals a recurring pattern in chemistry:
each new energy source or instrument reveals a new layer of reality.
\end{historybox}

\begin{puzzlebox}
Lavoisier listed ``caloric'' (heat) as an element,
reasoning that it could not be created or destroyed,
could pass between bodies, and behaved like a fluid.

\begin{enumerate}[label=(\alph*), leftmargin=*]
  \item In what ways does caloric theory \emph{fit} the
        observations of heat flow?
        (Hint: what happens when you put a hot and cold object together?)

  \item Benjamin Thompson (Count Rumford) drilled cannon barrels
        and noticed that the boring process generated heat
        \emph{indefinitely} as long as the horse turned the drill.
        Why is this a fatal problem for caloric theory?

  \item If heat were a fluid substance, what experimental result
        would you expect if you ground two identical pieces of metal
        together for a long time inside a sealed, insulated container?
        What did Rumford actually observe?
\end{enumerate}
\end{puzzlebox}

% ─── Section: The Law of Definite Proportions ─────────────────────────────────
\subsection{The Law of Definite Proportions}

With Lavoisier's framework established, other chemists began
systematically measuring the mass ratios in which elements combine.

The results were startling in their regularity.

\begin{center}
\begin{tabular}{llccc}
\toprule
\textbf{Compound} & \textbf{Formula} &
\textbf{Element 1 (mass \%)} &
\textbf{Element 2 (mass \%)} &
\textbf{Mass ratio} \\
\midrule
Water             & \ce{H2O}     & H: 11.1\%  & O: 88.9\%  & H:O = 1:8    \\
Carbon dioxide    & \ce{CO2}     & C: 27.3\%  & O: 72.7\%  & C:O = 3:8    \\
Ammonia           & \ce{NH3}     & N: 82.3\%  & H: 17.7\%  & N:H = 14:3   \\
Iron(II) sulfide  & \ce{FeS}     & Fe: 63.5\% & S: 36.5\%  & Fe:S = 7:4   \\
Magnesium oxide   & \ce{MgO}     & Mg: 60.3\% & O: 39.7\%  & Mg:O = 3:2   \\
\bottomrule
\end{tabular}
\end{center}

\vspace{0.6em}

The French chemist Joseph Louis Proust measured
these ratios obsessively between 1797 and 1804,
working with samples from different sources, different continents,
different methods of preparation.
The ratios never changed.

\begin{namedconcept}{Law of Definite Proportions (Proust, 1799)}
Every pure chemical compound contains its elements
in \emph{fixed, definite proportions by mass},
regardless of the source of the compound or the
method by which it was prepared.
\end{namedconcept}

\vspace{0.5em}

This law sounds dry.
Its implication is explosive.

\begin{puzzlebox}
Suppose matter is \emph{continuous} --- like a perfectly
smooth, infinitely divisible fluid.
You are mixing hydrogen and oxygen to make water.

\begin{enumerate}[label=(\alph*), leftmargin=*]
  \item If matter is continuous, is there any reason
        why hydrogen and oxygen should combine in
        a fixed ratio rather than any ratio you choose?

  \item Now suppose instead that hydrogen and oxygen are made
        of tiny indivisible particles (call them $h$ and $o$),
        and that water is formed when exactly \emph{one} $h$
        combines with exactly \emph{one} $o$.
        \begin{itemize}[nosep]
          \item Would you expect a fixed mass ratio? Why?
          \item What would the fixed ratio tell you about
                the relative masses of $h$ and $o$?
        \end{itemize}

  \item The actual mass ratio for water is $\text{H}:\text{O} = 1:8$.
        If water is \ce{HO} (one hydrogen atom per one oxygen atom),
        what is the mass of an oxygen atom relative to hydrogen?
        If water is \ce{H2O} (two hydrogens per one oxygen),
        what is the relative mass of oxygen?
        How would you decide which formula is correct?
\end{enumerate}
\end{puzzlebox}

\vspace{0.5em}

This puzzle is not rhetorical.
It is \emph{exactly} the reasoning that led John Dalton
to the atomic hypothesis.
The fixed ratios demand a discrete structure.

% ─── Section: The Law of Multiple Proportions ─────────────────────────────────
\subsection{A Second Clue: The Law of Multiple Proportions}

If the law of definite proportions was surprising,
the next discovery was even stranger.

Carbon and oxygen form \emph{two} different compounds.
In one (carbon monoxide, \ce{CO}), the ratio C:O by mass is
approximately 3:4.
In the other (carbon dioxide, \ce{CO2}), the ratio C:O by mass is
approximately 3:8.

Notice what happens when you fix the amount of carbon
and compare the oxygen in the two compounds:

\begin{center}
\begin{tikzpicture}[
  every node/.style={font=\sffamily\small},
  bar/.style={fill=ModuleColor!60, draw=ModuleColor, minimum width=0.8cm}
]
  % Carbon monoxide bar
  \node at (0, 2.5) {\textbf{Carbon monoxide}};
  \node[bar, minimum height=1.5cm] at (-0.5, 0.75) {};
  \node[bar, fill=DiscoveryColor!60, draw=DiscoveryColor,
        minimum height=2.0cm, minimum width=0.8cm] at (0.5, 1.0) {};
  \node[font=\sffamily\footnotesize, text=ModuleColor]
        at (-0.5, -0.2) {C: 3 g};
  \node[font=\sffamily\footnotesize, text=DiscoveryColor]
        at (0.5, -0.2) {O: 4 g};

  % Carbon dioxide bar
  \node at (5, 2.5) {\textbf{Carbon dioxide}};
  \node[bar, minimum height=1.5cm] at (4.5, 0.75) {};
  \node[bar, fill=DiscoveryColor!60, draw=DiscoveryColor,
        minimum height=4.0cm, minimum width=0.8cm] at (5.5, 2.0) {};
  \node[font=\sffamily\footnotesize, text=ModuleColor]
        at (4.5, -0.2) {C: 3 g};
  \node[font=\sffamily\footnotesize, text=DiscoveryColor]
        at (5.5, -0.2) {O: 8 g};

  % Ratio annotation
  \draw[{Stealth}-{Stealth}, thick, color=HistoryColor]
        (1.2, 1.0) -- (4.0, 1.0)
        node[midway, above, font=\sffamily\footnotesize,
             text=HistoryColor]{O ratio = 4:8 = \textbf{1:2}};
\end{tikzpicture}
\end{center}

\vspace{0.5em}

For a fixed mass of carbon, the masses of oxygen in the two
compounds are in the ratio \textbf{1:2} --- a small whole number.

Dalton noticed this for many pairs of compounds and stated it
as a general law.

\begin{namedconcept}{Law of Multiple Proportions (Dalton, 1803)}
When two elements combine to form more than one compound,
the masses of one element that combine with a fixed mass
of the other are in the ratio of \emph{small whole numbers}.
\end{namedconcept}

\vspace{0.5em}

Why \emph{small whole numbers}?
This is the decisive clue.
Continuous matter has no reason to produce integer ratios.
Only discrete particles do.

\begin{center}
\begin{tikzpicture}[
  every node/.style={font=\sffamily\small, align=center},
  atom/.style={circle, minimum size=0.65cm, font=\sffamily\tiny\bfseries,
               draw, thick}
]
  % CO: one C, one O
  \node[font=\sffamily\small\bfseries, color=JourneyColor]
        at (0, 1.5) {Carbon monoxide: \ce{CO}};
  \node[atom, fill=ModuleColor!30, draw=ModuleColor] (c1) at (-0.5, 0) {C};
  \node[atom, fill=DiscoveryColor!30, draw=DiscoveryColor] (o1) at (0.5, 0) {O};
  \draw[thick, JourneyColor!60] (c1) -- (o1);

  % CO2: one C, two O
  \node[font=\sffamily\small\bfseries, color=JourneyColor]
        at (6, 1.5) {Carbon dioxide: \ce{CO2}};
  \node[atom, fill=ModuleColor!30, draw=ModuleColor] (c2) at (5.5, 0) {C};
  \node[atom, fill=DiscoveryColor!30, draw=DiscoveryColor] (o2a) at (6.7, 0.5) {O};
  \node[atom, fill=DiscoveryColor!30, draw=DiscoveryColor] (o2b) at (6.7,-0.5) {O};
  \draw[thick, JourneyColor!60] (c2) -- (o2a);
  \draw[thick, JourneyColor!60] (c2) -- (o2b);

  % Explanation
  \node[font=\sffamily\footnotesize\itshape, text=HistoryColor,
        align=center]
        at (3, -1.2)
        {Each C gets 1 oxygen in CO, and 2 oxygens in CO$_2$.\\
         The ratio of oxygen per carbon is exactly 1:2.\\
         This is only possible if oxygen comes in discrete units.};
\end{tikzpicture}
\end{center}

\begin{problembox}
\begin{enumerate}[label=\textbf{\arabic*.}]
  \item Nitrogen and oxygen form several compounds.
        In three of them, the mass of oxygen that combines
        with \SI{14.0}{\gram} of nitrogen is:
        \begin{center}
        \begin{tabular}{lll}
        \toprule
        Compound & Formula & O mass per \SI{14}{\gram} N \\
        \midrule
        Nitrous oxide   & \ce{N2O}  & \SI{8.0}{\gram}  \\
        Nitric oxide    & \ce{NO}   & \SI{16.0}{\gram} \\
        Nitrogen dioxide & \ce{NO2} & \SI{32.0}{\gram} \\
        \bottomrule
        \end{tabular}
        \end{center}
        \begin{enumerate}[label=(\alph*)]
          \item Show that these masses are in the ratio 1:2:4.
          \item Explain why this ratio is consistent with
                the law of multiple proportions.
          \item If matter were continuous,
                could you have a compound with O:N = 1.73:1
                (an irrational number)?
                Would atoms allow this?
        \end{enumerate}

  \item Sulfur and oxygen form two compounds:
        sulfur dioxide (\ce{SO2}) and sulfur trioxide (\ce{SO3}).
        In \ce{SO2}, the mass ratio S:O is approximately 1:1.
        \begin{enumerate}[label=(\alph*)]
          \item Predict the mass ratio S:O in \ce{SO3}.
          \item Verify your prediction:
                sulfur has atomic mass \SI{32}{\gram\per\mol},
                oxygen \SI{16}{\gram\per\mol}.
          \item What is the ratio of oxygen in \ce{SO2} to
                oxygen in \ce{SO3}
                (for the same amount of sulfur)?
        \end{enumerate}
\end{enumerate}
\end{problembox}

% ─── Section: Dalton's Atomic Theory ──────────────────────────────────────────
\subsection{Dalton's Atomic Theory: The Inevitable Conclusion}

By 1803, John Dalton --- a quiet English schoolteacher with
a passion for meteorology and gas behaviour ---
had assembled the evidence into a coherent theory.
He was not the first to suggest atoms existed.
He was the first to make the idea \emph{quantitative}
and connect it directly to experimental data.

\begin{namedconcept}{Dalton's Atomic Theory (1803)}
\begin{enumerate}[leftmargin=2em, nosep]
  \item All matter is composed of extremely small, indivisible
        particles called \newterm{atoms}.
  \item Atoms of a given element are identical in mass
        and properties.
  \item Atoms of different elements have different masses.
  \item Compounds are formed when atoms of different elements
        combine in fixed, small whole-number ratios.
  \item Chemical reactions are the rearrangement of atoms.
        No atoms are created, destroyed, or changed into
        atoms of a different element.
\end{enumerate}
\end{namedconcept}

\vspace{0.5em}

Notice what this theory accomplishes in one stroke:

\begin{center}
\begin{tikzpicture}[
  every node/.style={font=\sffamily\small, align=center},
  arrow/.style={-{Stealth[length=5pt]}, thick, color=ModuleColor}
]
  % Laws on the left
  \node[draw=JourneyColor!50, fill=JourneyColor!8, rounded corners=3pt,
        minimum width=4.0cm, minimum height=0.8cm]
        (cons) at (0, 2.0) {Conservation of mass};
  \node[draw=JourneyColor!50, fill=JourneyColor!8, rounded corners=3pt,
        minimum width=4.0cm, minimum height=0.8cm]
        (def)  at (0, 0.8) {Definite proportions};
  \node[draw=JourneyColor!50, fill=JourneyColor!8, rounded corners=3pt,
        minimum width=4.0cm, minimum height=0.8cm]
        (mult) at (0,-0.4) {Multiple proportions};

  % Dalton's theory in center
  \node[draw=DiscoveryColor, fill=DiscoveryColor!10, rounded corners=5pt,
        minimum width=3.5cm, minimum height=1.4cm,
        font=\sffamily\small\bfseries]
        (dalton) at (5.5, 0.8) {Dalton's\\Atomic Theory};

  \foreach \src in {cons, def, mult}
    \draw[arrow] (\src.east) -- (dalton.west);

  % What it predicts on right
  \node[draw=ModuleColor!50, fill=ModuleColor!8, rounded corners=3pt,
        minimum width=4.0cm, minimum height=0.8cm]
        (p1) at (11, 2.0) {Explains conservation of mass};
  \node[draw=ModuleColor!50, fill=ModuleColor!8, rounded corners=3pt,
        minimum width=4.0cm, minimum height=0.8cm]
        (p2) at (11, 0.8) {Explains definite proportions};
  \node[draw=ModuleColor!50, fill=ModuleColor!8, rounded corners=3pt,
        minimum width=4.0cm, minimum height=0.8cm]
        (p3) at (11,-0.4) {Explains multiple proportions};
  \node[draw=ModuleColor!50, fill=ModuleColor!8, rounded corners=3pt,
        minimum width=4.0cm, minimum height=0.8cm]
        (p4) at (11,-1.6) {\emph{Predicts} atomic masses};

  \foreach \tgt in {p1, p2, p3, p4}
    \draw[arrow] (dalton.east) -- (\tgt.west);
\end{tikzpicture}
\end{center}

\vspace{0.5em}

Three laws become consequences of one underlying picture.
And the theory makes a new prediction:
we should be able to determine the \emph{relative masses}
of atoms from the mass ratios in compounds.

\begin{historybox}
\textbf{Dalton's Atomic Symbols}

Before Berzelius introduced the modern letter-based notation,
Dalton drew elaborate circular symbols for each element.
Hydrogen was a plain circle,
oxygen a circle with a dot,
nitrogen a circle with a vertical line, and so on.
Water was depicted as a circle-with-dot touching a plain circle.

Dalton assumed water was \ce{HO} (one atom of each).
This gave him the wrong atomic mass for oxygen
(he calculated 8, when it is actually 16).
He did not know the formula was \ce{H2O}.
That would be corrected by Gay-Lussac's law of combining volumes
and Avogadro's hypothesis --- a story we will pick up in Journey~2.
\end{historybox}

% ─── Section: Relative Atomic Masses ──────────────────────────────────────────
\subsection{Relative Atomic Masses: Weighing the Invisible}

Dalton could not see atoms.
But he could \emph{weigh} them --- relatively.

The logic is elegant.
If water is \ce{H2O} and the mass ratio H:O is 1:8,
then one oxygen atom must be 16 times heavier than
one hydrogen atom.
(Because 2 hydrogen atoms weigh 1 part, and 1 oxygen atom
weighs 8 parts, so oxygen : hydrogen = 8 : (1/2) = 16 : 1.)

\begin{center}
\begin{tikzpicture}[
  every node/.style={font=\sffamily\small},
  atom/.style={circle, draw, thick, font=\sffamily\footnotesize\bfseries}
]
  % Water molecule
  \node[atom, minimum size=1.6cm,
        fill=DiscoveryColor!25, draw=DiscoveryColor]
        (O) at (0,0) {O, 16 u};
  \node[atom, minimum size=0.9cm,
        fill=blue!20, draw=blue!70]
        (H1) at (-1.6, 0.8) {H, 1 u};
  \node[atom, minimum size=0.9cm,
        fill=blue!20, draw=blue!70]
        (H2) at (-1.6,-0.8) {H, 1 u};
  \draw[thick, JourneyColor!60] (O) -- (H1);
  \draw[thick, JourneyColor!60] (O) -- (H2);

  % Mass annotation
  \node[font=\sffamily\footnotesize, text=JourneyColor!70]
        at (2.5, 0.8) {Total mass = 18 u};
  \node[font=\sffamily\footnotesize, text=JourneyColor!70]
        at (2.5, 0.2) {H fraction = 2/18 = 11.1\%};
  \node[font=\sffamily\footnotesize, text=JourneyColor!70]
        at (2.5,-0.4) {O fraction = 16/18 = 88.9\%};
  \node[font=\sffamily\footnotesize\itshape, text=HistoryColor]
        at (2.5,-1.0) {Matches Lavoisier's measurement!};
\end{tikzpicture}
\end{center}

\vspace{0.5em}

The unit ``u'' stands for \newterm{atomic mass unit},
defined so that hydrogen has mass 1.
This is purely a relative scale; we are not yet
measuring actual masses in grams.

Berzelius spent years measuring mass ratios for dozens of compounds
to build the first table of relative atomic masses.
His values for many elements were accurate to within 1\%.

\subsubsection{The Mole: Connecting Atoms to Grams}

The atomic mass unit is useful for relative comparisons,
but chemists need to work with laboratory quantities.
The bridge is the \newterm{mole}.

\begin{namedconcept}{The Mole}
One \newterm{mole} of a substance contains
exactly $6.022 \times 10^{23}$ particles
(Avogadro's number, $N_A$).

The \newterm{molar mass} of a substance is the mass
of one mole in grams, numerically equal to the
relative atomic (or molecular) mass.

\begin{center}
$n = \dfrac{m}{M}$
\end{center}

where $n$ is moles, $m$ is mass in grams,
and $M$ is molar mass in \si{\gram\per\mol}.
\end{namedconcept}

\vspace{0.5em}

The mole is simply a convenient counting unit,
like a ``dozen'' for eggs.
We use it because atoms are so small that even a tiny quantity
of substance contains an incomprehensibly large number of them.

\begin{center}
\begin{tikzpicture}[
  every node/.style={font=\sffamily\small, align=center}
]
  % Scale comparison
  \node[draw=JourneyColor!50, fill=JourneyColor!8, rounded corners=3pt,
        minimum width=5cm, minimum height=1.4cm]
        (lab) at (0,0)
        {\textbf{Laboratory scale}\\
         \SI{18}{\gram} of water = 1 mole\\
         contains $6.022\times10^{23}$ molecules};

  \node[draw=ModuleColor!50, fill=ModuleColor!8, rounded corners=3pt,
        minimum width=5cm, minimum height=1.4cm]
        (atom) at (7.5,0)
        {\textbf{Atomic scale}\\
         1 water molecule\\
         mass = $3.0\times10^{-23}\ \si{\gram}$};

  \draw[-{Stealth[length=5pt]}, thick, color=DiscoveryColor]
        (lab.east) -- (atom.west)
        node[midway, above, font=\sffamily\footnotesize]{$\times N_A$};
  \draw[-{Stealth[length=5pt]}, thick, color=DiscoveryColor]
        (atom.west) -- (lab.east)
        node[midway, below, font=\sffamily\footnotesize]{$\div N_A$};
\end{tikzpicture}
\end{center}

\begin{problembox}
\begin{enumerate}[label=\textbf{\arabic*.}]
  \item Calculate the molar mass of each compound.
        Then find the number of moles in the given sample.
        \begin{enumerate}[label=(\alph*)]
          \item \ce{H2O}: \SI{36.0}{\gram}
          \item \ce{CO2}: \SI{22.0}{\gram}
          \item \ce{NaCl}: \SI{58.5}{\gram}
          \item \ce{CaCO3}: \SI{50.0}{\gram}
        \end{enumerate}

  \item \SI{4.0}{\gram} of hydrogen gas (\ce{H2}) reacts
        completely with oxygen to form water.
        \begin{enumerate}[label=(\alph*)]
          \item How many moles of \ce{H2} is this?
          \item Using the balanced equation
                \ce{2H2 + O2 -> 2H2O},
                how many moles of \ce{O2} are needed?
          \item What mass of \ce{O2} is required?
          \item What mass of water is produced?
        \end{enumerate}

  \item An unknown compound contains only carbon and hydrogen.
        Complete combustion of \SI{2.40}{\gram} of the compound
        produces \SI{7.92}{\gram} of \ce{CO2} and
        \SI{3.24}{\gram} of \ce{H2O}.
        \begin{enumerate}[label=(\alph*)]
          \item Find the mass of carbon in the sample.
                (Hint: all carbon ends up in \ce{CO2}.)
          \item Find the mass of hydrogen in the sample.
          \item Find the simplest whole-number ratio of
                C atoms to H atoms.
          \item What is the \emph{empirical formula} of
                the compound?
        \end{enumerate}
\end{enumerate}
\end{problembox}

% ─── Section: Empirical and Molecular Formulas ────────────────────────────────
\subsection{Empirical Formulas vs Molecular Formulas}

Dalton's theory gives us the \emph{ratio} of atoms in a compound,
but not always the actual number.

\begin{namedconcept}{Empirical Formula vs Molecular Formula}
The \newterm{empirical formula} gives the simplest whole-number
ratio of atoms in a compound.

The \newterm{molecular formula} gives the actual number of
each atom in one molecule.
It is always a whole-number multiple of the empirical formula.
\end{namedconcept}

\vspace{0.5em}

\begin{center}
\begin{tabular}{llll}
\toprule
\textbf{Compound} & \textbf{Empirical} & \textbf{Molecular} & \textbf{Multiple} \\
\midrule
Water             & \ce{H2O}    & \ce{H2O}   & $\times 1$ \\
Hydrogen peroxide & \ce{HO}     & \ce{H2O2}  & $\times 2$ \\
Glucose           & \ce{CH2O}   & \ce{C6H12O6} & $\times 6$ \\
Benzene           & \ce{CH}     & \ce{C6H6}  & $\times 6$ \\
Ethane            & \ce{CH3}    & \ce{C2H6}  & $\times 2$ \\
Acetylene         & \ce{CH}     & \ce{C2H2}  & $\times 2$ \\
\bottomrule
\end{tabular}
\end{center}

\vspace{0.5em}

Notice that benzene and acetylene have the \emph{same empirical
formula} (\ce{CH}) but are completely different compounds
with completely different properties.
To determine the molecular formula, we need additional information:
the \emph{molar mass} of the compound.

\begin{puzzlebox}
Glucose (\ce{C6H12O6}) and formaldehyde (\ce{CH2O})
have the same empirical formula.
Glucose is the primary fuel for the brain.
Formaldehyde is a toxic preservative used for embalming.

\begin{enumerate}[label=(\alph*), leftmargin=*]
  \item What are the molar masses of glucose and formaldehyde?
  \item By what factor do the molecular formulas differ?
  \item What does this example illustrate about the relationship
        between empirical formula and chemical properties?
  \item We will return to this question in Journey~6.
        Write your best current hypothesis for why the
        same empirical formula can produce such different substances.
\end{enumerate}
\end{puzzlebox}

% ─── Section: Where Dalton's Theory Breaks Down ───────────────────────────────
\subsection{Where Dalton's Theory Breaks Down}

Dalton's atomic theory was a triumph.
But every theory has a boundary,
and great scientists know that finding the boundary
is as important as establishing the theory.

\begin{modelbreaks}
Dalton's theory, as originally stated, makes several claims
that turned out to be wrong or incomplete.

\textbf{1. Atoms are not indivisible.}\\
In 1897, J.~J.~Thomson discovered that atoms contain
\emph{electrons} --- particles far smaller than atoms
that can be removed from them.
Atoms have internal structure.

\textbf{2. Atoms of the same element are not always identical.}\\
In the 20th century, \newterm{isotopes} were discovered:
atoms of the same element with different masses
(different numbers of neutrons in the nucleus).
Chlorine, for instance, exists naturally as a mixture of
\ce{^{35}Cl} (mass 35) and \ce{^{37}Cl} (mass 37),
which is why its average atomic mass is 35.5.

\textbf{3. Atoms \emph{can} be converted to other elements.}\\
In nuclear reactions, one element transforms into another.
Dalton's statement that atoms cannot change into a different
element holds only for \emph{chemical} reactions.

\textbf{4. The assumption of simplest formulas was wrong.}\\
Dalton assumed water was \ce{HO}.
It is \ce{H2O}.
Without independent evidence (such as gas volumes),
the atomic mass scale is ambiguous.

\vspace{0.4em}
\emph{Despite these failures, Dalton's core insight ---
that matter is made of discrete particles whose ratios
determine composition --- remains completely correct.
The failures refined the theory; they did not overturn it.}
\end{modelbreaks}

% ─── Section: Deep Problems ───────────────────────────────────────────────────
\subsection*{Deep Problems}

\begin{deepbox}
\begin{enumerate}[label=\textbf{D\arabic*.}]

  \item \textbf{The Avogadro Number from First Principles.}\\
        The actual numerical value of Avogadro's number
        ($N_A = 6.022\times10^{23}$) could not be determined
        from Dalton's theory alone.
        Dalton only established \emph{relative} atomic masses.
        To find absolute masses, we need an independent measurement.

        \begin{enumerate}[label=(\alph*)]
          \item Why can we not deduce $N_A$ from mass ratios alone?
                What additional type of measurement is needed?

          \item One method uses Brownian motion:
                Robert Brown observed in 1827 that pollen grains
                suspended in water jiggle randomly.
                Einstein showed in 1905 that the statistics of
                this jiggling could be used to calculate $N_A$.
                Qualitatively, what does the jiggling tell us
                about the structure of the water?

          \item Millikan's oil drop experiment (1909) measured
                the charge of a single electron to be
                $e = 1.602\times10^{-19}\ \si{\coulomb}$.
                The Faraday constant (charge to deposit 1 mole
                of a singly-charged ion in electrolysis) is
                $F = \SI{96485}{\coulomb\per\mol}$.
                Use these to calculate $N_A$.
        \end{enumerate}

  \item \textbf{Reasoning About Isotopes.}\\
        Chlorine has two stable isotopes:
        \ce{^{35}Cl} (mass exactly 35 u, natural abundance 75.77\%)
        and \ce{^{37}Cl} (mass exactly 37 u, natural abundance 24.23\%).

        \begin{enumerate}[label=(\alph*)]
          \item Calculate the weighted average atomic mass of chlorine.
                Verify it matches the periodic table value of 35.45 u.

          \item If Dalton had used natural chlorine in his experiments,
                why would the mass ratios he measured still obey
                the law of definite proportions?
                (Hint: what is the probability distribution of
                \ce{^{35}Cl} vs \ce{^{37}Cl} in a large sample?)

          \item Dalton's postulate says atoms of the same element are
                identical. Isotopes violate this.
                But the law of multiple proportions still holds.
                Explain why the existence of isotopes does not
                destroy Dalton's framework,
                even though it contradicts his postulate.
        \end{enumerate}

  \item \textbf{Building the Atom Count.}\\
        One mole of water (\SI{18}{\gram}) contains
        $6.022\times10^{23}$ molecules.
        Each molecule contains 3 atoms.
        The diameter of a water molecule is approximately
        $3.2\times10^{-10}\ \si{\metre}$.

        \begin{enumerate}[label=(\alph*)]
          \item How many atoms are in \SI{18}{\gram} of water?

          \item If you lined up all the water molecules in
                one mole in a single row,
                how long would the line be?
                Express your answer in light-years
                ($1\ \text{ly} \approx 9.46\times10^{15}\ \si{\metre}$).

          \item The volume of one water molecule is approximately
                $(3.2\times10^{-10})^3\ \si{\metre\cubed}$.
                The volume of \SI{18}{\gram} of liquid water is
                \SI{18}{\milli\litre}.
                Estimate the packing fraction
                (fraction of the volume actually occupied by molecules).
                What does this tell you about the space \emph{between}
                molecules?
        \end{enumerate}

\end{enumerate}
\end{deepbox}

% ─── Reflection ───────────────────────────────────────────────────────────────
\subsection*{Reflection}

\begin{reflectionbox}
\begin{enumerate}[label=\textbf{R\arabic*.}]
  \item Dalton's atomic theory explains three macroscopic laws
        (conservation of mass, definite proportions, multiple proportions)
        in terms of one microscopic picture.
        \textbf{What is the logical structure of this kind of argument?}
        Is Dalton \emph{proving} that atoms exist,
        or is he proposing the \emph{simplest hypothesis consistent}
        with the data?
        What would it take to \emph{prove} atoms exist?

  \item The law of multiple proportions gives ratios of small whole numbers:
        1:2, 1:3, 2:3, and so on.
        \textbf{Why must these be \emph{small} integers,
        not large ones like 1:247?}
        What does this tell us about chemical bonding?
        (You do not need to answer this fully ---
        write your current hypothesis.)

  \item We found that the empirical formula alone does not
        determine the molecular formula.
        Glucose and formaldehyde have the same empirical formula
        but utterly different properties.
        \textbf{What additional structural feature must be responsible
        for the difference?}
        Is composition the whole story?

  \item Dalton assumed water was \ce{HO} and got the wrong atomic
        mass for oxygen.
        \textbf{What does this tell us about the relationship
        between a theory and its predictions?}
        A theory can be deeply right in its core structure
        but wrong in details.
        How do scientists decide which parts of a wrong prediction
        to fix and which parts to keep?
\end{enumerate}
\end{reflectionbox}

\vspace{1.5em}
\begin{center}
\begin{tcolorbox}[
  enhanced, colback=DiscoveryColor!8, colframe=DiscoveryColor,
  width=0.88\textwidth, arc=3pt
]
\centering\sffamily\small
\textbf{What we invented in this module.}\\[4pt]
We invented the atom --- not by seeing it,
but by being \emph{forced to} by the data.
The fixed mass ratios in compounds are not just a curiosity.
They are a logical demand: only a world made of discrete,
indivisible particles with fixed relative masses
can produce such ratios.
Dalton's genius was not in proposing that atoms might exist.
It was in showing that the existing experimental measurements
\emph{require} their existence.\\[4pt]
This pattern --- experimental regularity demanding a hidden structure ---
will appear again and again as we go deeper into chemistry.
The hydrogen spectrum will demand quantised energy levels.
The reactivity of the noble gases will demand electron shells.
Each new regularity is a door that can only be opened
by proposing something invisible.
\end{tcolorbox}
\end{center}

\vspace{2em}
\noindent\textcolor{JourneyColor!40}{\rule{\textwidth}{0.6pt}}
\begin{center}
\textit{\sffamily\small
Module 1.3 complete. Continue to Module 1.4:
\textbf{Lavoisier's Conservation and the Quantitative Revolution}---
or proceed to Journey~2:
\textbf{Architecture of the Atom}.}
\end{center}

\end{document}